\section{Accounting Context}
\label{sec:related}

\subsection{A heterogeneous inventory sector}

IPCC Volume~3 organises IPPU into mineral industries (2.A), chemical industries (2.B),
metal industries (2.C), non-energy products from fuels and solvents (2.D), electronics
(2.E), fluorinated substitutes for ozone-depleting substances (2.F), and other product
manufacture and use (2.G) \citep{ipcc2006ippu}. The 2019 Refinement updates important methods
and factors but does not remove the central data requirement: the activity must correspond to
the physical process that emits.

This is why headline production quantities can be misleading. Cement output is not clinker
production. Total fertiliser output differs from ammonia production. Billets do not identify
the steelmaking route or its carbon inputs, and annual HFC leakage depends on the applications
and equipment stocks in which the chemicals are used. Each quantity is valuable evidence, but
the relevant process definition is also required.

\subsection{Material balances as shared accounting}

Several IPPU categories become clearer when written as conserved material flows. For cement,
domestic clinker production plus imports minus exports and stock accumulation must equal the
clinker ground into cement. For ammonia, feedstock carbon becomes gross process \COtwo{}, which
is divided among atmospheric emissions, urea production, other recovery and any permanent
storage. For fluorinated gases, annual supply enters equipment banks and leaves through
manufacturing losses, operating leakage, retirement, recovery and destruction.

These balances complement IPCC emission factors by constraining the activity bases to
which those factors apply and expose residuals that a collection of independent equations can
hide. They also connect IPPU to neighbouring sectors without double counting: cement kiln
fuel remains Energy 1.A.2, while clinker calcination remains IPPU 2.A.1.

\subsection{Evaluation and official reporting}

Broad global products such as EDGAR are useful only after year, category, gas and GWP are
matched. A national ``Processes'' total cannot validate cement alone. Pakistan's own CRT is a
stronger accounting comparator because it exposes category rows, native gases, activity data,
implied factors, recovery and notation keys \citep{pakcrt2026}. We therefore use the CRT as an
official method audit and scope comparator. Public production, trade and consumption datasets
provide separate evidence for activity magnitude and reporting boundaries.
